\section{Related Work}\label{sec:related}
% Budget: 1 page. Draft. The own ECSA 2025 paper (Okumus2025) is cited in third person.

\paragraph{Contract-Based Design.} Benveniste et al.~\cite{Benveniste2018} develop the meta-theory of contracts, including refinement and composition, on which Sect.~\ref{sec:contracts} builds. Cimatti and Tonetta~\cite{CimattiTonetta2015} decompose the contracts of an architecture along its component hierarchy and reduce the correctness of contract refinement to proof obligations in temporal logic that are discharged at design time. These works establish refinement; we ask what an established refinement implies for monitors that observe finite prefixes of a behaviour.

\paragraph{Monitor Synthesis from Contracts and Requirements.} Okumus et al.~\cite{Okumus2025} propose a taxonomy of monitoring aspects for service-oriented automotive architectures and synthesise monitors from contracts attached to service interfaces; correctness of the monitors and relationships between components are not addressed. Guissouma et al.~\cite{Guissouma2023} monitor contracts of evolving learning-based components, and Watanabe et al.~\cite{Watanabe2018} monitor signal temporal logic specifications within contract-based design. Perez et al.~\cite{Perez2022} translate structured natural-language requirements into monitors of the Copilot framework~\cite{Pike2013}. Sokolsky et al.~\cite{Sokolsky2016} verify guarantees before deployment and monitor the assumptions of a component afterwards, and Xu et al.~\cite{Xu2019} decompose a root contract on non-functional properties into a hierarchy of monitored contracts. Closest to our work, Luckcuck et al.~\cite{Luckcuck2022} specify the nodes of robotic systems by assume/guarantee contracts in first-order logic, derive system-level properties from them by inference rules, and translate the contracts into runtime monitors. None of these works proves the monitors correct with respect to the contract semantics or states under which condition the silence of component monitors implies a system-level property on a finite observation.
% TODO: read the full text of Luckcuck2022; check whether the derivation of system-level properties is stated for finite runs and whether monitor correctness is argued. The sentence above rests on abstracts only. Bauer and Falcone~\cite{BauerFalcone2016} decentralise the monitoring of a global temporal formula by distributing sub-formulae among local monitors that exchange obligations. In contrast, our decomposition is given by the contracts of the architecture and justified by refinement; the component monitors do not communicate.

\paragraph{Safety and Finite Observations.} That a violated safety property has a finite bad prefix is the basis of runtime verification~\cite{KupfermanVardi2001}, and monitors for linear temporal logic are constructed to reject every bad prefix as early as possible~\cite{BauerLeuckerSchallhart2011}. Cimatti et al.~\cite{Cimatti2024} exploit assumptions on the system to obtain earlier verdicts. Conditions that relate finite to infinite behaviours are also known from the composition of specifications: Abadi and Lamport~\cite{AbadiLamport1995} require that every finite behaviour satisfying the safety part of a specification can be extended to one satisfying the whole specification. Extendability (Def.~\ref{def:extendable}) is a condition of this kind. Our contribution is not the notion but its role: it is exactly what is needed to transfer a contract refinement to component monitors that do not look ahead, and Theorem~\ref{thm:deadend} reduces it to a check on finite prefixes.

\paragraph{Verified Monitors.} Schneider et al.~\cite{Schneider2019} and Basin et al.~\cite{Basin2020} verify a monitor for metric first-order temporal logic in Isabelle/HOL; their logic is considerably more expressive than ours, in particular through quantification over data. Finkbeiner et al.~\cite{Finkbeiner2020} compile stream-based specifications to Rust code with annotations from which the correctness of each generated monitor is proved automatically, and Scott et al.~\cite{Scott2023} validate each compiled Copilot monitor against its specification by a bisimulation argument. Mitsch and Platzer~\cite{MitschPlatzer2016} derive monitors by proof from verified hybrid system models, such that offline verification results carry over to the monitored execution. These works assure a single monitor. We verify the synthesis once and, beyond that, mechanise the compositional layer that relates the monitors of an architecture to each other.

\paragraph{Monitoring of Learning-Based Perception.} Ferreira et al.~\cite{Ferreira2025} survey safety monitors for perception functions and classify them by what they observe: inputs, internal representations, outputs, or external sensors. These mechanisms detect that a prediction should not be trusted; they are complementary to our work, as their verdicts or scores enter contracts as atoms. Astorga et al.~\cite{Astorga2023} introduce perception contracts that bound the error of a perception component such that a system-level invariant is preserved, and synthesise them from data; such contracts are candidates for the guarantees our monitors check at runtime.

\paragraph{Blame.} Attributing a contract violation to one of the parties originates from contracts for higher-order functions~\cite{FindlerFelleisen2002}. Theorem~\ref{thm:blame} transfers the idea to components communicating over streams.
